% Please do not change %
\documentclass[11pt]{article} % font size
\usepackage[margin=1in]{geometry} % margins
\usepackage{amsmath,amsthm,amssymb} % for the  common "math symbols"
\usepackage{algorithm}
\usepackage{algpseudocode}
\usepackage[shortlabels]{enumitem} % for enumeration
\usepackage{booktabs} % if you need tables
\usepackage{listings}
\usepackage{color}
\usepackage{quiver}
\usepackage{ulem}
\DeclareMathOperator*{\argmax}{argmax}
\DeclareMathOperator*{\argmin}{argmin}

% \theoremstyle{definition}
\newtheorem{theorem}{Theorem}[section]
\newtheorem{corollary}[theorem]{Corollary}
\newtheorem{lemma}[theorem]{Lemma}
\newtheorem{proposition}[theorem]{Proposition}
\newtheorem{claim}[theorem]{Claim}
\newtheorem{conjecture}[theorem]{Conjecture}
\theoremstyle{definition}
\newtheorem{definition}[theorem]{Definition}
\theoremstyle{definition}
\newtheorem{fact}[theorem]{Fact}
\theoremstyle{definition}
\newtheorem{example}[theorem]{Example}
\newtheorem{note}[theorem]{Note}

\lstset{
  frame=none,
  xleftmargin=2pt,
  stepnumber=1,
  numbers=left,
  numbersep=5pt,
  numberstyle=\ttfamily\tiny\color[gray]{0.3},
  belowcaptionskip=\bigskipamount,
  captionpos=b,
  escapeinside={*'}{'*},
  language=haskell,
  tabsize=2,
  emphstyle={\bf},
  commentstyle=\it,
  stringstyle=\mdseries\rmfamily,
  showspaces=false,
  keywordstyle=\bfseries\rmfamily,
  columns=flexible,
  basicstyle=\small\sffamily,
  showstringspaces=false,
  morecomment=[l]\%,
}
% Feel free to include additional packages (if needed), but make sure it does not override the margin/font requirements.
%
\setlength{\parindent}{0pt} % no indent for new paragraphs.
\setlength{\parskip}{5pt} % distance between paragraphs, which will be used when you have a blank line in your LaTeX code.


%%%%%%%%%%%%%%%%%%%%%%%%%%%%%%%%%%%%%%%%%%%%%%%%%%%%%%%%%%%%%%%%%%%%%%%%%%%%%%
\title{Categories, Proofs, and Processes: \\Chapter 3 Limits}
%
\author{Wira Azmoon Ahmad}
\date{1 Feb 2024}
%
\begin{document}
%
\maketitle

%You may use the usual \LaTeX\ environments to structure your answers:
%\begin{align}
%        \label{basegnn}
%        \mathbf{h}_u^{(t+1)} = \sigma \bigg( \mathbf{W}_\text{self}^{(t)} \mathbf{h}_u^{(t)}  + \mathbf{W}_\text{neigh}^{(t)} \sum_{v \in N(u)} \mathbf{h}_v^{(t)} + \mathbf{b}^{(t)} \bigg): 
%\end{align}


%\begin{table}[h]
%\caption{A simple table.}
%\centering
%\begin{tabular}[t]{lrrrr}
%\toprule
%Graphs & Can & Be  & Fun & Too \\ 
%\midrule 
%$G_1$ & $1$  & $2$ & $3$ & $4$ \\ 
%$G_2$ & $-1$  & $-2$ & $-3$ & $-4$ \\ 
%$G_3$ & $0.1$  & $0.2$ & $0.3$ & $0.4$ \\ 
%\bottomrule
%\end{tabular}
%\label{tab}
%\end{table}

% A functor $F:\mathcal{C} \rightarrow \mathcal{D}$ is
% \begin{itemize}
%     \item \textbf{faithful} if each map $F_{A,B}:\mathcal{C}(A,B) \rightarrow \mathcal{D}(F\ A, F\ B)$ is injective;
%     \item \textbf{full} if each map $F_{A,B}:\mathcal{C}(A,B) \rightarrow \mathcal{D}(F\ A, F\ B)$ is surjective
% \end{itemize}


\addtocounter{section}{2}
\section{Limits and Colimits}
\subsection{Equalisers}
\subsubsection{Motivation}
In \textbf{Set}, we can define equality between functions $f,g$ when $\forall x\in X, f(x)=g(x)$,
but we might be interested in situations where where $f(x)=g(x)$ for $x\in Y\subsetneq X$.

For example, $X=\{0,1\}^{*}$, but only a subset $G$ could encode graphs using this alphabet.

\subsubsection{Definition}
\begin{definition}
Consider a pair of parallel arrows $A\overset{f}{\underset{g}{\rightrightarrows}}B$.
An \textbf{\textit{\underline{equaliser}}} of $(f,g)$ is an arrow $e: E\rightarrow A$ such that
\[f\circ e = g\circ e\]
and for any $h:D\rightarrow A$ such that $f\circ h = g\circ h$, 
there is a unique $\hat{h}: D\rightarrow E$ so that
\[h = e\circ \hat{h}\]
Diagrammatically,
% https://q.uiver.app/#q=WzAsNCxbMSwwLCJBIl0sWzIsMCwiQiJdLFswLDAsIkUiXSxbMCwxLCJEIl0sWzAsMSwiZiIsMCx7Im9mZnNldCI6LTF9XSxbMCwxLCJnIiwyLHsib2Zmc2V0IjoxfV0sWzIsMCwiZSJdLFszLDAsImgiLDJdLFszLDIsIlxcaGF0e2h9IiwwLHsic3R5bGUiOnsiYm9keSI6eyJuYW1lIjoiZGFzaGVkIn19fV1d
\[\begin{tikzcd}
	E & A & B \\
	D
	\arrow["f", shift left, from=1-2, to=1-3]
	\arrow["g"', shift right, from=1-2, to=1-3]
	\arrow["e", from=1-1, to=1-2]
	\arrow["h"', from=2-1, to=1-2]
	\arrow["{\hat{h}}", dashed, from=2-1, to=1-1]
\end{tikzcd}\]
\end{definition}

\begin{note}
    A diagram \underline{commutes} if any two paths with a common source and target, 
    with at least one of the paths at least length $>1$, are equal.
\end{note}

\begin{example}
    In \textbf{Set}, the equaliser of $f,g$ is given by the inclusion
    \[\{x\in A | f(x)=g(x)\} \hookrightarrow A\]
    For $f,g: \mathbb{R} \rightarrow \mathbb{R}$ where 
    \[\forall x, f(x)=ax^2+bx+c, g(x) =0\]
    the equaliser is the roots of $f(x)$ if they exist.
\end{example}

\subsubsection{Results}
\begin{proposition}
    Any equaliser in any category is monic.
\end{proposition}
\begin{proof}
% https://q.uiver.app/#q=WzAsNCxbMSwwLCJBIl0sWzIsMCwiQiJdLFswLDAsIkUiXSxbMCwxLCJEIl0sWzAsMSwiZiIsMCx7Im9mZnNldCI6LTF9XSxbMCwxLCJnIiwyLHsib2Zmc2V0IjoxfV0sWzIsMCwiZSJdLFszLDAsImgiLDJdLFszLDIsIngiLDAseyJvZmZzZXQiOi0xfV0sWzMsMiwieSIsMix7Im9mZnNldCI6MX1dXQ==
\[\begin{tikzcd}
	E & A & B \\
	D
	\arrow["f", shift left, from=1-2, to=1-3]
	\arrow["g"', shift right, from=1-2, to=1-3]
	\arrow["e", from=1-1, to=1-2]
	\arrow["h"', from=2-1, to=1-2]
	\arrow["x", shift left, from=2-1, to=1-1]
	\arrow["y"', shift right, from=2-1, to=1-1]
\end{tikzcd}\]
Suppose $e \circ x = e\circ y$, then it is equal to $h$. Then
\[f\circ h = f\circ e\circ x = g\circ e\circ x = g\circ h\]
There is a unique map $h' :D\rightarrow E$ such that $e\circ h' = h$.
But since
\[e\circ x = h \wedge e\circ y = h \implies x = y = h'\]
by uniqueness, so $e$ is monic.
\end{proof}

\begin{definition}
    \textbf{Coequaliser} in $\mathcal{C}$ is an equaliser in $\mathcal{C}^{op}$.
\end{definition}
\begin{proposition}
    Any coequaliser in any category is epic.
\end{proposition}

\begin{proposition}
    Equalisers and coequalisers are unique up to a unique isomorphism.
\end{proposition}
\begin{proof}
Suppose
% https://q.uiver.app/#q=WzAsNCxbMSwwLCJBIl0sWzIsMCwiQiJdLFswLDAsIkVfMSJdLFswLDEsIkVfMiJdLFswLDEsImYiLDAseyJvZmZzZXQiOi0xfV0sWzAsMSwiZyIsMix7Im9mZnNldCI6MX1dLFsyLDAsImVfMSJdLFszLDAsImVfMiIsMl1d
\[\begin{tikzcd}
	{E_1} & A & B \\
	{E_2}
	\arrow["f", shift left, from=1-2, to=1-3]
	\arrow["g"', shift right, from=1-2, to=1-3]
	\arrow["{e_1}", from=1-1, to=1-2]
	\arrow["{e_2}"', from=2-1, to=1-2]
\end{tikzcd}\]
Thus there are unique maps
\[u_1: E_2 \rightarrow E_1, \quad e_2 = e_1 \circ u_1\]
\[u_2: E_1 \rightarrow E_2, \quad e_1 = e_2 \circ u_2\]
Then identities exist
\[1_{E_1}: E_1 \rightarrow E_1 := u_1 \circ u_2\]
\[1_{E_2}: E_2 \rightarrow E_2 := u_2 \circ u_1\]
\end{proof}

\subsection{Pullbacks}
\subsubsection{Definitions}
\begin{definition}
    Given a pair of morphisms
    \[A\xrightarrow[]{f} C \xleftarrow[]{g}B\]
    (called a cospan),
    we define an $(f,g)$-cone to be 
    \[A\xleftarrow[]{p}D \xrightarrow[]{q}B\]
    such that the following diagram commutes.
    % https://q.uiver.app/#q=WzAsNCxbMCwwLCJEIl0sWzEsMCwiQiJdLFswLDEsIkEiXSxbMSwxLCJDIl0sWzAsMSwicCJdLFswLDIsInAiLDJdLFsyLDMsImYiLDJdLFsxLDMsImciXV0=
\[\begin{tikzcd}
	D & B \\
	A & C
	\arrow["q", from=1-1, to=1-2]
	\arrow["p"', from=1-1, to=2-1]
	\arrow["f"', from=2-1, to=2-2]
	\arrow["g", from=1-2, to=2-2]
\end{tikzcd}\]
\end{definition}
\begin{definition}
    The \textbf{\textit{\underline{pullback}}}\underline{ of f along g} is an $(f,g)$-cone
    such that for any other $(f,g)$-cone $A\xleftarrow[]{p}D \xrightarrow[]{q}B$, 
    there is a unique $h:D'\rightarrow D$ such that $p'=p\circ h, q'=q\circ h$.
    Diagrammatically,
    % https://q.uiver.app/#q=WzAsNSxbMSwxLCJEIl0sWzIsMSwiQiJdLFsxLDIsIkEiXSxbMiwyLCJDIl0sWzAsMCwiRCciXSxbMCwxLCJxIl0sWzAsMiwicCIsMl0sWzIsMywiZiIsMl0sWzEsMywiZyJdLFs0LDEsInEnIiwwLHsiY3VydmUiOi0yfV0sWzQsMiwicCciLDIseyJjdXJ2ZSI6Mn1dLFs0LDAsImgiLDAseyJzdHlsZSI6eyJib2R5Ijp7Im5hbWUiOiJkYXNoZWQifX19XV0=
\[\begin{tikzcd}
	{D'} \\
	& D & B \\
	& A & C
	\arrow["q", from=2-2, to=2-3]
	\arrow["p"', from=2-2, to=3-2]
	\arrow["f"', from=3-2, to=3-3]
	\arrow["g", from=2-3, to=3-3]
	\arrow["{q'}", curve={height=-12pt}, from=1-1, to=2-3]
	\arrow["{p'}"', curve={height=12pt}, from=1-1, to=3-2]
	\arrow["h", dashed, from=1-1, to=2-2]
\end{tikzcd}\]
\end{definition}
\begin{note}
    Pullback often denoted as 
    \[A \times_C B\]
\end{note}
\begin{example}
    In \textbf{Set}, pullback is subset of cartesian product.
    Take a function $f:A\rightarrow B$ and subset $V\subseteq B$. Consider
    \[f^{-1}(V) = \{a\in A | f(a)\in V\} \subseteq A\]
    Consider
    % https://q.uiver.app/#q=WzAsNCxbMCwwLCJmXnstMX0oVikiXSxbMSwwLCJWIl0sWzAsMSwiQSJdLFsxLDEsIkIiXSxbMCwxLCJcXGJhcntmfSJdLFswLDIsIiIsMix7InN0eWxlIjp7InRhaWwiOnsibmFtZSI6Imhvb2siLCJzaWRlIjoidG9wIn19fV0sWzIsMywiZiIsMl0sWzEsMywiIiwwLHsic3R5bGUiOnsidGFpbCI6eyJuYW1lIjoiaG9vayIsInNpZGUiOiJ0b3AifX19XV0=
\[\begin{tikzcd}
	{f^{-1}(V)} & V \\
	A & B
	\arrow["{\bar{f}}", from=1-1, to=1-2]
	\arrow[hook, from=1-1, to=2-1]
	\arrow["f"', from=2-1, to=2-2]
	\arrow[hook, from=1-2, to=2-2]
\end{tikzcd}\]
So $f^{-1}(V)$ is a pullback, unique up to a unique isomorphism.
\end{example}
\begin{proposition}
    Pullbacks are unique up to a unique isomorphism.
\end{proposition}

\subsubsection{Results}
\begin{proposition}
    If a category $\mathcal{C}$ has binary products and equalisers, then it has pullbacks.
\end{proposition}
\begin{proof}
    Given the cospan
    \[A\xleftarrow[]{p}C \xrightarrow[]{q}B\]
    we want to construct its pullback using products and equalisers.
    Consider the product 
    \[A\xleftarrow[]{\pi_1}A\times B \xrightarrow[]{\pi_2}B\]
    Form an equaliser $e:E \rightarrow A\times B$.
    The following illustrative diagram does not commute
    % https://q.uiver.app/#q=WzAsNSxbMSwxLCJBXFx0aW1lcyBCIl0sWzIsMSwiQiJdLFsxLDIsIkEiXSxbMiwyLCJDIl0sWzAsMCwiRSJdLFswLDEsIlxccGlfMiJdLFswLDIsIlxccGlfMSIsMl0sWzIsMywiZiIsMl0sWzEsMywiZyJdLFs0LDEsIlxccGlfMiBcXGNpcmMgZSIsMCx7ImN1cnZlIjotMn1dLFs0LDIsIlxccGlfMSBcXGNpcmMgZSIsMix7ImN1cnZlIjoyfV0sWzQsMCwiZSJdXQ==
\[\begin{tikzcd}
	E \\
	& {A\times B} & B \\
	& A & C
	\arrow["{\pi_2}", from=2-2, to=2-3]
	\arrow["{\pi_1}"', from=2-2, to=3-2]
	\arrow["f"', from=3-2, to=3-3]
	\arrow["g", from=2-3, to=3-3]
	\arrow["{\pi_2 \circ e}", curve={height=-12pt}, from=1-1, to=2-3]
	\arrow["{\pi_1 \circ e}"', curve={height=12pt}, from=1-1, to=3-2]
	\arrow["e", from=1-1, to=2-2]
\end{tikzcd}\]
We want to show $A\xleftarrow[]{\pi_1 \circ e}E \xrightarrow[]{\pi_2 \circ e}B$ is a pullback,
since $e$ might not be unique.
Assume the following diagrams commute
% https://q.uiver.app/#q=WzAsOCxbMCwwLCJEIl0sWzEsMCwiQiJdLFswLDEsIkEiXSxbMSwxLCJDIl0sWzQsMCwiRCJdLFs0LDIsIkEiXSxbNiwwLCJCIl0sWzUsMSwiQVxcdGltZXMgQiJdLFswLDEsImgiXSxbMCwyLCJrIiwyXSxbMiwzLCJmIiwyXSxbMSwzLCJnIl0sWzQsNSwiayIsMl0sWzQsNiwiaCJdLFs0LDcsIlxcbGFuZ2xlIGssIGhcXHJhbmdsZSIsMV0sWzcsNiwiXFxwaV8yIiwyXSxbNyw1LCJcXHBpXzEiXV0=
\[\begin{tikzcd}
	D & B &&& D && B \\
	A & C &&&& {A\times B} \\
	&&&& A
	\arrow["h", from=1-1, to=1-2]
	\arrow["k"', from=1-1, to=2-1]
	\arrow["f"', from=2-1, to=2-2]
	\arrow["g", from=1-2, to=2-2]
	\arrow["k"', from=1-5, to=3-5]
	\arrow["h", from=1-5, to=1-7]
	\arrow["{\langle k, h\rangle}"{description}, from=1-5, to=2-6]
	\arrow["{\pi_2}"', from=2-6, to=1-7]
	\arrow["{\pi_1}", from=2-6, to=3-5]
\end{tikzcd}\]
Then 
\[f\circ \pi_1 \circ \langle k, h\rangle = g\circ \pi_2 \langle k, h\rangle\]
Combine this with the original equaliser and we have
% https://q.uiver.app/#q=WzAsNCxbMCwwLCJFIl0sWzEsMCwiQVxcdGltZXMgQiJdLFsyLDAsIkMiXSxbMCwxLCJEIl0sWzAsMV0sWzEsMiwiZ1xcY2lyY1xccGlfMiIsMix7Im9mZnNldCI6MX1dLFsxLDIsImZcXGNpcmMgXFxwaV8xIiwwLHsib2Zmc2V0IjotMX1dLFszLDAsInUiLDAseyJzdHlsZSI6eyJib2R5Ijp7Im5hbWUiOiJkYXNoZWQifX19XSxbMywxLCJcXGxhbmdsZSBrLGhcXHJhbmdsZSIsMl1d
\[\begin{tikzcd}
	E & {A\times B} & C \\
	D
	\arrow[from=1-1, to=1-2]
	\arrow["{g\circ\pi_2}"', shift right, from=1-2, to=1-3]
	\arrow["{f\circ \pi_1}", shift left, from=1-2, to=1-3]
	\arrow["u", dashed, from=2-1, to=1-1]
	\arrow["{\langle k,h\rangle}"', from=2-1, to=1-2]
\end{tikzcd}\]
Unique map $u$, thus $A\xleftarrow[]{\pi_1 \circ e}E \xrightarrow[]{\pi_2 \circ e}B$ is a pullback.
\end{proof}

\subsection{Summary}
Constructs follow similar pattern
\begin{enumerate}
    \item Determine shape of a diagram
    \item Highlight an object and arrows in that diagram
    \item Stipulate unique map to (from) that object from (to) all other objects with the same 
    "composition" properties
\end{enumerate}
\begin{note}
    We are exhibiting a \underline{universal (mapping) property} where the unique map is a "unique morphism".
\end{note}

\subsection{Limits}


\end{document}
